\chapter{Implementation}
\label{ch:5}


\begin{figure}[H]
\centering
\includegraphics[width=0.8\textwidth]{04_fl_sequence.png}
\caption{Federated Learning Training Sequence}
\label{fig:04_fl_sequence}
\end{figure}


\section{5.1 Development Environment Setup}

\subsection{5.1.1 Infrastructure Requirements}
Hardware:

\begin{itemize}
  \item GPU: NVIDIA Tesla P100/T4 (16GB+ VRAM) recommended
  \item RAM: 32GB system memory (for full dataset preloading)
  \item Storage: 100GB+ (raw data + models + results)
  \item Network: ≥ 10Mbps (for dataset downloads)
\end{itemize}

Software Stack:

\begin{itemize}
  \item Python 3.10+
  \item PyTorch 2.0+
  \item MONAI 1.3+ (medical imaging framework)
  \item scikit-learn 1.3+
  \item NumPy 1.24+, Pandas
  \item Matplotlib/Seaborn (visualization)
  \item Kaggle API (dataset access)
\end{itemize}

Recommended Platforms:

\begin{itemize}
  \item Kaggle Notebooks (free GPU, simplified setup)
  \item Google Colab (free T4 GPU)
  \item Local machine with NVIDIA GPU + CUDA 11.8/12.1
  \item Docker container (reproducible environment)
\end{itemize}


\begin{figure}[H]
\centering
\includegraphics[width=0.85\textwidth]{07_progression_pipeline.png}
\caption{Progression Module Workflow}
\label{fig:07_progression_pipeline}
\end{figure}


\section{5.2 Module 1: Federated Classification Implementation}

\subsection{5.2.1 Model Definition: ResNet-18 Brain Tumor Classifier}
File: federated\_learning/src/model.py

"""
ResNet-18 model for brain tumor classification.
Pretrained on ImageNet; final FC layer replaced for 3-class output.
"""
import torch
import torch.nn as nn
import torchvision.models as models

class BrainTumorResNet(nn.Module):
    """
    ResNet-18 with custom FC head for brain tumor classification.
    Input: (B, 3, 224, 224) RGB images
    Output: (B, num\_classes) logits
    """
    def \_\_init\_\_(self, num\_classes=3, pretrained=True):
        super().\_\_init\_\_()
        \# Load pretrained ResNet-18
        weights = models.ResNet18\_Weights.DEFAULT if pretrained else None
        self.model = models.resnet18(weights=weights)
        
        \# Replace final FC layer
        num\_features = self.model.fc.in\_features  \# 512
        self.model.fc = nn.Linear(num\_features, num\_classes)
    
    def forward(self, x):
        """Forward pass"""
        return self.model(x)

def create\_model(num\_classes=3, device="cuda"):
    """Instantiate model, move to device, print summary"""
    model = BrainTumorResNet(num\_classes=num\_classes, pretrained=True)
    model = model.to(device)
    
    total\_params = sum(p.numel() for p in model.parameters())
    trainable\_params = sum(p.numel() for p in model.parameters() if p.requires\_grad)
    mb = total\_params * 4 / (1024 ** 2)
    
    print(f"Model: ResNet-18 | Params: \{total\_params:,\} | "
          f"Trainable: \{trainable\_params:,\} | Size: \~\{mb:.1f\}MB")
    return model


\subsection{5.2.2 Federated Client Implementation}
File: federated\_learning/src/client.py

"""
FederatedClient: Handles local training for one hospital/client.
"""
import copy
import torch
import torch.nn as nn
import torch.optim as optim
from tqdm.notebook import tqdm

class FederatedClient:
    """
    Represents one hospital/silo in federated setting.
    Typical round:
        client.set\_model(global\_model)
        client.set\_optimizer(lr=0.001)
        weights, losses, accs = client.train\_local(epochs=5)
        val\_loss, val\_acc = client.validate()
    """
    def \_\_init\_\_(self, client\_id, train\_loader, val\_loader, device="cuda"):
        self.client\_id = client\_id
        self.train\_loader = train\_loader
        self.val\_loader = val\_loader
        self.device = device
        self.model = None
        self.optimizer = None
        self.criterion = nn.CrossEntropyLoss()
        self.dataset\_size = len(train\_loader.dataset)
    
    def set\_model(self, global\_model):
        """Deep-copy server's global model to local client"""
        self.model = copy.deepcopy(global\_model)
        self.model.to(self.device)
    
    def set\_optimizer(self, learning\_rate=0.001):
        """Initialize optimizer with client's model parameters"""
        self.optimizer = optim.Adam(self.model.parameters(), lr=learning\_rate)
    
    def train\_local(self, epochs=5, verbose=False):
        """
        Local training on private hospital data.
        Returns: (state\_dict, train\_losses, train\_accs)
        """
        self.model.train()
        epoch\_losses, epoch\_accs = [], []
        
        for ep in range(epochs):
            running\_loss = 0.0
            correct = total = 0
            
            pbar = tqdm(self.train\_loader, 
                       desc=f"\{self.client\_id\} Epoch \{ep+1\}/\{epochs\}",
                       disable=not verbose, leave=False)
            
            for images, labels in pbar:
                images = images.to(self.device)
                labels = labels.to(self.device)
                
                self.optimizer.zero\_grad()
                outputs = self.model(images)
                loss = self.criterion(outputs, labels)
                loss.backward()
                self.optimizer.step()
                
                running\_loss += loss.item()
                \_, predicted = outputs.max(1)
                total += labels.size(0)
                correct += predicted.eq(labels).sum().item()
            
            ep\_loss = running\_loss / len(self.train\_loader)
            ep\_acc = 100.0 * correct / total
            epoch\_losses.append(ep\_loss)
            epoch\_accs.append(ep\_acc)
        
        return self.model.state\_dict(), epoch\_losses, epoch\_accs
    
    def validate(self):
        """Evaluate on local validation set. Returns (loss, accuracy\%)"""
        self.model.eval()
        val\_loss = 0.0
        correct = total = 0
        
        with torch.no\_grad():
            for images, labels in self.val\_loader:
                images = images.to(self.device)
                labels = labels.to(self.device)
                
                outputs = self.model(images)
                loss = self.criterion(outputs, labels)
                
                val\_loss += loss.item()
                \_, predicted = outputs.max(1)
                total += labels.size(0)
                correct += predicted.eq(labels).sum().item()
        
        val\_loss /= len(self.val\_loader)
        val\_acc = 100.0 * correct / total
        return val\_loss, val\_acc


\subsection{5.2.3 QPSO Server Implementation (Core Aggregation)}
File: federated\_learning/src/server\_qpso.py (excerpt)

"""
QPSOServer: Quantum Particle Swarm Optimization for FL aggregation.
Maintains pbest (personal best), gbest (global best), mbest (mean best).
"""
import copy
import torch
import math

class QPSOServer:
    """Central server using QPSO-based aggregation strategy."""
    
    def \_\_init\_\_(self, global\_model, clients, device="cuda", beta=0.7):
        self.global\_model = global\_model
        self.clients = clients
        self.device = device
        self.beta = beta
        
        \# Per-client personal bests and scores
        self.personal\_best = \{\}
        self.personal\_best\_scores = \{\}
        
        \# Global best
        self.global\_best = None
        self.global\_best\_score = 0.0
        
        \# Mean best (computed each round)
        self.mean\_best = None
        print(f"QPSO Server initialized | Clients=\{len(clients)\} | β=\{beta\}")
    
    def initialize\_particles(self):
        """Set all pbest and gbest to current global model"""
        state = copy.deepcopy(self.global\_model.state\_dict())
        for c in self.clients:
            self.personal\_best[c.client\_id] = copy.deepcopy(state)
            self.personal\_best\_scores[c.client\_id] = 0.0
        self.global\_best = copy.deepcopy(state)
        self.global\_best\_score = 0.0
        print("✅ QPSO particles initialized")
    
    def update\_personal\_best(self, client\_id, weights, val\_acc):
        """Update client's personal best if new accuracy is better"""
        if val\_acc > self.personal\_best\_scores[client\_id]:
            self.personal\_best[client\_id] = copy.deepcopy(weights)
            self.personal\_best\_scores[client\_id] = val\_acc
            return True
        return False
    
    def update\_global\_best(self, client\_id, val\_acc):
        """Update global best if this client is better"""
        if val\_acc > self.global\_best\_score:
            self.global\_best = copy.deepcopy(self.personal\_best[client\_id])
            self.global\_best\_score = val\_acc
            return True
        return False
    
    def calculate\_mean\_best(self):
        """Compute element-wise mean of all personal-best state\_dicts"""
        first\_id = self.clients[0].client\_id
        self.mean\_best = copy.deepcopy(self.personal\_best[first\_id])
        
        \# Average all pbests
        for key in self.mean\_best:
            sum\_tensor = torch.zeros\_like(self.mean\_best[key])
            for client in self.clients:
                sum\_tensor += self.personal\_best[client.client\_id][key]
            self.mean\_best[key] = sum\_tensor / len(self.clients)
    
    def aggregate\_qpso(self):
        """
        QPSO aggregation step:
        For each parameter: x\_new = p ± β * |mbest - x| * ln(1/u)
        where p = φ*mbest + (1-φ)*gbest is the attraction point
        """
        self.calculate\_mean\_best()
        w\_new = copy.deepcopy(self.global\_best)
        
        for param\_name in w\_new:
            \# Per-parameter quantum update
            phi = torch.rand(1).item()  \# ∈ [0,1]
            u = torch.rand(1).item()  \# ∈ [0,1], then clamp to [0.3, 1.0]
            u = max(u, 0.3)
            
            \# Attraction point
            p = (phi * self.mean\_best[param\_name] + 
                 (1 - phi) * self.global\_best[param\_name])
            
            \# Quantum position update
            if u > 0:
                ln\_term = math.log(1.0 / u)
                perturbation = (self.beta * 
                               torch.abs(self.mean\_best[param\_name] - 
                                        self.global\_best[param\_name]) * 
                               ln\_term)
                
                \# Clamp perturbation for stability
                perturbation = torch.clamp(perturbation, -0.1, 0.1)
                
                \# Random sign
                sign = 1 if torch.rand(1).item() > 0.5 else -1
                w\_new[param\_name] = p + sign * perturbation
        
        \# Load aggregated weights into global model
        self.global\_model.load\_state\_dict(w\_new)
        return self.global\_best\_score
    
    def federated\_round(self, round\_num, test\_loader):
        """Execute one complete FL round"""
        \# Broadcast global model to all clients
        for client in self.clients:
            client.set\_model(self.global\_model)
        
        \# Local training
        client\_weights = \{\}
        for client in self.clients:
            \_, \_, \_ = client.train\_local(epochs=5, verbose=False)
            val\_loss, val\_acc = client.validate()
            
            client\_weights[client.client\_id] = client.model.state\_dict()
            self.update\_personal\_best(client.client\_id, 
                                     client.model.state\_dict(), 
                                     val\_acc)
            self.update\_global\_best(client.client\_id, val\_acc)
        
        \# QPSO Aggregation
        best\_score = self.aggregate\_qpso()
        
        \# Evaluate on global test set
        test\_acc = self.evaluate\_on\_test\_set(test\_loader)
        
        return best\_score, test\_acc
    
    def evaluate\_on\_test\_set(self, test\_loader):
        """Evaluate global model on balanced test set"""
        self.global\_model.eval()
        correct = total = 0
        
        with torch.no\_grad():
            for images, labels in test\_loader:
                images = images.to(self.device)
                labels = labels.to(self.device)
                outputs = self.global\_model(images)
                \_, predicted = outputs.max(1)
                total += labels.size(0)
                correct += predicted.eq(labels).sum().item()
        
        return 100.0 * correct / total


\section{5.3 Module 2: Segmentation Implementation}

\subsection{5.3.1 Data Loading and Preprocessing (MONAI)}
"""
Segmentation preprocessing using MONAI framework.
Handles 3D MRI volumes with 4 modalities (T1, T1ce, T2, FLAIR).
"""
from monai.transforms import (
    Compose, LoadImaged, NormalizeIntensityd, 
    Orientationd, Spacingd, RandCropByPosNegLabel,
    EnsureChannelFirstd, EnsureTyped
)
from monai.data import Dataset, DataLoader

def get\_preprocessing\_transforms(training=True):
    """Define preprocessing pipeline for BraTS data"""
    transforms = Compose([
        LoadImaged(keys=["image", "label"]),
        EnsureChannelFirstd(keys=["image", "label"]),
        EnsureTyped(keys=["image", "label"]),
        
        \# Reorient to RAS standard
        Orientationd(keys=["image", "label"], axcodes="RAS"),
        
        \# Resample to 1mm³ isotropic
        Spacingd(keys=["image", "label"], 
                pixdim=(1.0, 1.0, 1.0),
                mode=("bilinear", "nearest")),
        
        \# Normalize intensity (nonzero voxels only)
        NormalizeIntensityd(keys=["image"], 
                           nonzero=True, 
                           channel\_wise=True),
        
        \# Random crop during training (96³ patches, balanced pos/neg)
        RandCropByPosNegLabel(keys=["image", "label"],
                             label\_key="label",
                             spatial\_size=(96, 96, 96),
                             pos=1.0,  \# 50\% positive (tumor) samples
                             neg=1.0,  \# 50\% negative (background)
                             num\_samples=4) if training else None,
    ])
    
    \# Filter None
    return Compose([t for t in transforms if t is not None])

\# Load dataset
train\_transforms = get\_preprocessing\_transforms(training=True)
train\_ds = Dataset(data=train\_files, transform=train\_transforms)
train\_loader = DataLoader(train\_ds, batch\_size=1, shuffle=True, num\_workers=2)


\subsection{5.3.2 Inference: Sliding Window with Model}
"""
3D Segmentation inference using sliding window approach.
Handles large volumes by processing overlapping patches.
"""
from monai.inferers import sliding\_window\_inference

def segment\_patient(model, patient\_image\_path, device):
    """
    Segment one patient's 3D MRI volume.
    Args:
        model: Trained AttentionUnet
        patient\_image\_path: Path to .nii.gz file
        device: torch device (cuda/cpu)
    Returns:
        predictions: (1, 3, H, W, D) segmentation mask
    """
    \# Load and preprocess
    sample = train\_transforms(\{"image": patient\_image\_path, 
                              "label": patient\_image\_path\})
    image = sample["image"].unsqueeze(0).to(device)  \# (1, 4, H, W, D)
    
    model.eval()
    with torch.no\_grad():
        \# Sliding window inference: process in 96³ patches
        outputs = sliding\_window\_inference(
            inputs=image,
            roi\_size=(96, 96, 96),
            sw\_batch\_size=4,
            predictor=model,
            overlap=0.5  \# 50\% overlap between patches
        )
        
        \# Apply sigmoid + threshold
        predictions = (outputs.sigmoid() > 0.5).float()
    
    return predictions  \# (1, 3, H, W, D)

\# Example usage
model = AttentionUnet(spatial\_dims=3, in\_channels=4, out\_channels=3)
model.load\_state\_dict(torch.load("best\_model.pth", map\_location=device))

pred = segment\_patient(model, "patient\_T1.nii.gz", device)
print(f"Segmentation shape: \{pred.shape\}")
print(f"Tumor voxels: TC=\{pred[0,0].sum():.0f\}, "
      f"WT=\{pred[0,1].sum():.0f\}, ET=\{pred[0,2].sum():.0f\}")


\section{5.4 Module 3: Progression Implementation}

\subsection{5.4.1 Mathematical Model Fitting}
File: progression/src/01\_mathematical\_models.py (excerpt)

"""
Fit mathematical models to tumor volume time series.
"""
import numpy as np
from scipy.optimize import curve\_fit

class MathematicalModels:
    """Classical tumor growth models"""
    
    @staticmethod
    def logistic(t, V0, K, r):
        """V(t) = K / (1 + ((K-V0)/V0)*exp(-r*t))"""
        return K / (1 + ((K - V0) / V0) * np.exp(-r * t))
    
    @staticmethod
    def gompertz(t, V0, a, b):
        """V(t) = V0 * exp((a/b)*(1-exp(-b*t)))"""
        return V0 * np.exp((a / b) * (1 - np.exp(-b * t)))
    
    @staticmethod
    def exponential(t, V0, lam):
        """V(t) = V0 * exp(λ*t)"""
        return V0 * np.exp(lam * t)

def fit\_models\_per\_patient(times, volumes, grade="HGG"):
    """
    Fit all models to patient data.
    Args:
        times: Array of time points (days from first scan)
        volumes: Array of tumor volumes (mm³)
        grade: 'HGG' or 'LGG' for model parameter initialization
    Returns:
        best\_model\_name, best\_r2, fitted\_params
    """
    results = \{\}
    
    \# Logistic fit
    try:
        V0 = volumes[0]
        K = volumes.max() * 1.5
        r = 0.01
        popt, \_ = curve\_fit(MathematicalModels.logistic, 
                           times, volumes,
                           p0=[V0, K, r],
                           bounds=([V0*0.5, K, 0.0001], 
                                   [V0*1.5, K*2, 1.0]))
        pred = MathematicalModels.logistic(times, *popt)
        r2 = 1 - np.sum((volumes - pred)**2) / np.sum((volumes - volumes.mean())**2)
        results['logistic'] = \{'params': popt, 'r2': r2, 'pred': pred\}
    except:
        pass
    
    \# Select best model
    best\_model = max(results, key=lambda k: results[k]['r2'])
    return best\_model, results[best\_model]['r2'], results[best\_model]['params']


\subsection{5.4.2 LSTM Hybrid Training}
File: progression/src/07\_hybrid\_lstm\_training.py (excerpt)

"""
Train LSTM to learn residuals from mathematical model.
"""
import torch
import torch.nn as nn

class ResidualLSTM(nn.Module):
    """LSTM learns where mathematical model fails"""
    
    def \_\_init\_\_(self, input\_size=1, hidden\_size=32, num\_layers=1):
        super().\_\_init\_\_()
        self.lstm = nn.LSTM(input\_size, hidden\_size, num\_layers, batch\_first=True)
        self.fc1 = nn.Linear(hidden\_size, 64)
        self.fc2 = nn.Linear(64, 32)
        self.fc3 = nn.Linear(32, 1)
        self.relu = nn.ReLU()
    
    def forward(self, x):
        """x: (batch, seq\_len, 1) residual sequence"""
        lstm\_out, (h\_n, c\_n) = self.lstm(x)
        last\_hidden = lstm\_out[:, -1, :]  \# (batch, hidden\_size)
        
        x = self.relu(self.fc1(last\_hidden))
        x = self.relu(self.fc2(x))
        residual\_correction = self.fc3(x)
        
        return residual\_correction

def train\_hybrid\_lstm(model, train\_loader, val\_loader, device, epochs=50):
    """Train LSTM on residual sequences"""
    optimizer = torch.optim.Adam(model.parameters(), lr=0.001)
    criterion = nn.MSELoss()
    
    for epoch in range(epochs):
        model.train()
        train\_loss = 0.0
        
        for x\_batch, y\_batch in train\_loader:
            x\_batch = x\_batch.to(device)
            y\_batch = y\_batch.to(device)
            
            optimizer.zero\_grad()
            y\_pred = model(x\_batch)
            loss = criterion(y\_pred, y\_batch)
            loss.backward()
            optimizer.step()
            
            train\_loss += loss.item()
        
        train\_loss /= len(train\_loader)
        
        \# Validation
        model.eval()
        val\_loss = 0.0
        with torch.no\_grad():
            for x\_batch, y\_batch in val\_loader:
                x\_batch = x\_batch.to(device)
                y\_batch = y\_batch.to(device)
                y\_pred = model(x\_batch)
                loss = criterion(y\_pred, y\_batch)
                val\_loss += loss.item()
        
        val\_loss /= len(val\_loader)
        
        if epoch \% 10 == 0:
            print(f"Epoch \{epoch\}: Train Loss=\{train\_loss:.6f\}, Val Loss=\{val\_loss:.6f\}")


\section{5.5 Key Implementation Details}

\subsection{5.5.1 Data Privacy in FL}
How FL achieves privacy:

\# PRIVATE (on Hospital Client)
hospital\_data = load\_private\_patient\_data()  \# Raw images
hospital\_model.train(hospital\_data)  \# Train locally
updated\_weights = hospital\_model.state\_dict()  \# Extract weights only

\# TRANSMITTED (Hospital → Server)
\# Only send: updated\_weights  (\~45 MB)
\# Never send: hospital\_data, intermediate activations, gradients

\# SERVER receives only weights, aggregates them
global\_weights = aggregate([hospital1\_weights, hospital2\_weights, hospital3\_weights])

\# Result: Hospital data never leaves premises ✓


\subsection{5.5.2 Federated Round Pseudocode}
\# Main FL Training Loop
for round\_t in range(1, 101):
    \# 1. Server broadcasts to all clients
    for client in clients:
        client.receive\_global\_model(global\_model)
    
    \# 2. Clients train locally (parallel)
    client\_weights = []
    for client in clients:
        weights, losses, accs = client.train\_local(epochs=5)
        val\_loss, val\_acc = client.validate()
        client\_weights.append((weights, val\_acc))
    
    \# 3. Server aggregates
    if strategy == "FedAvg":
        global\_weights = fedavg\_aggregate(client\_weights)
    elif strategy == "QPSO":
        global\_weights = qpso\_aggregate(client\_weights)
    
    \# 4. Evaluate on global test set
    test\_acc = evaluate(global\_model, test\_loader)
    
    \# Log metrics
    log(round=round\_t, test\_acc=test\_acc, fairness\_sigma=fairness\_metric)

\# Save final model
torch.save(global\_model.state\_dict(), "final\_fl\_model.pth")


\section{5.6 Performance Optimizations}

\subsection{5.6.1 GPU Memory Optimization}
\begin{itemize}
  \item **Gradient Checkpointing:** Save only select activations, recompute others during backprop
  \item **Mixed Precision:** Use float16 for forward pass, float32 for loss
  \item **Batch Size Tuning:** Max batch size before OOM; found to be 32 for ResNet-18 on 16GB GPU
  \item **Data Prefetching:** Load next batch while GPU processes current batch
\end{itemize}


\subsection{5.6.2 Communication Efficiency}
\begin{itemize}
  \item **Model Quantization (Future):** Compress weights from float32 → int8 (4x reduction)
  \item **Gradient Sparsification (Future):** Send only top-k gradients
  \item **Synchronous SGD (Current):** All clients must finish before aggregation (simple but suboptimal)
\end{itemize}


\section{5.7 Chapter Summary}
Figure: Federated Learning Sequence Diagram

Figure: LSTM Progression Forecasting Architecture

Figure: Progression Pipeline Workflow


\section{Chapter Summary}
This chapter described implementation details of all three modules using Python, PyTorch, MONAI, and SciPy with key code components.


\section{Key Code Components}

\subsection{BrainTumorResNet --- Classification Model}
\begin{lstlisting}[caption={ResNet-18 Brain Tumor Classifier (model.py)}, label={lst:model}]
"""
ResNet-18 model for 3-class brain tumour classification.
Pretrained on ImageNet; final FC layer replaced.
"""

import torch
import torch.nn as nn
import torchvision.models as models


class BrainTumorResNet(nn.Module):
    """
    ResNet-18 with the last FC layer replaced:
        512 → num_classes (default 3).
    """

    def __init__(self, num_classes=3, pretrained=True):
        super().__init__()
        weights = models.ResNet18_Weights.DEFAULT if pretrained else None
        self.model = models.resnet18(weights=weights)
        num_features = self.model.fc.in_features
        self.model.fc = nn.Linear(num_features, num_classes)

    def forward(self, x):
        return self.model(x)


def create_model(num_classes=3, device="cuda"):
    """Instantiate, move to device, and print a quick param summary."""
    model = BrainTumorResNet(num_classes=num_classes, pretrained=True)
    model = model.to(device)

    total  = sum(p.numel() for p in model.parameters())
    train  = sum(p.numel() for p in model.parameters() if p.requires_grad)
    mb     = total * 4 / (1024 ** 2)

    print(f"Model: ResNet-18  |  Params: {total:,} (trainable {train:,})  |  ~{mb:.1f} MB")
    return model
\end{lstlisting}

\subsection{FederatedClient --- Local Training}
\begin{lstlisting}[caption={Federated Client (client.py)}, label={lst:client}]
"""
FederatedClient — handles local training and validation for one FL client.
"""

import copy
import torch
import torch.nn as nn
import torch.optim as optim

try:
    from tqdm.notebook import tqdm
except ImportError:
    from tqdm import tqdm


class FederatedClient:
    """
    Represents one hospital / data silo in the federation.

    Typical round:
        client.set_model(global_model)
        client.set_optimizer(lr)
        weights, losses, accs = client.train_local(epochs)
        val_loss, val_acc = client.validate()
    """

    def __init__(self, client_id, train_loader, val_loader, device="cuda"):
        self.client_id    = client_id
        self.train_loader = train_loader
        self.val_loader   = val_loader
        self.device       = device
        self.model        = None
        self.optimizer    = None
        self.criterion    = nn.CrossEntropyLoss()
        self.dataset_size = len(train_loader.dataset)

    # ----- model / optimiser --------------------------------------------------

    def set_model(self, global_model):
        """Deep-copy the server's global model."""
        self.model = copy.deepcopy(global_model)
        self.model.to(self.device)

    def set_optimizer(self, learning_rate=0.001):
        self.optimizer = optim.Adam(self.model.parameters(), lr=learning_rate)

    # ----- local training -----------------------------------------------------

    def train_local(self, epochs=5, verbose=False):
        """
        Run local training on client data.

        Returns
        -------
        state_dict, list[float], list[float]
        """
        self.model.train()
        epoch_losses, epoch_accs = [], []

        for ep in range(epochs):
            running_loss = 0.0
            correct = total = 0

            pbar = tqdm(self.train_loader,
                        desc=f"{self.client_id} Epoch {ep+1}/{epochs}",
                        disable=not verbose, leave=False)

            for batch_idx, (images, labels) in enumerate(pbar):
                images, labels = images.to(self.device), labels.to(self.device)

                self.optimizer.zero_grad()
                outputs = self.model(images)
                loss = self.criterion(outputs, labels)
                loss.backward()
                self.optimizer.step()

                running_loss += loss.item()
                _, predicted  = outputs.max(1)
                total   += labels.size(0)
                correct += predicted.eq(labels).sum().item()

                if verbose:
                    pbar.set_postfix(Loss=f"{running_loss/(batch_idx+1):.4f}",
                                    Acc=f"{100.*correct/total:.2f}%")

            ep_loss = running_loss / len(self.train_loader)
            ep_acc  = 100.0 * correct / total
            epoch_losses.append(ep_loss)
            epoch_accs.append(ep_acc)

            if verbose:
                print(f"  {self.client_id} E{ep+1}: Loss={ep_loss:.4f}  Acc={ep_acc:.2f}%")

        return self.model.state_dict(), epoch_losses, epoch_accs

    # ----- validation ---------------------------------------------------------

    def validate(self):
        """Evaluate on the local validation set. Returns (loss, accuracy%)."""
        self.model.eval()
        val_loss = 0.0
        correct = total = 0

        with torch.no_grad():
            for images, labels in self.val_loader:
                images, labels = images.to(self.device), labels.to(self.device)
                outputs = self.model(images)
                loss    = self.criterion(outputs, labels)

                val_loss += loss.item()
                _, predicted = outputs.max(1)
                total   += labels.size(0)
                correct += predicted.eq(labels).sum().item()

        return val_loss / len(self.val_loader), 100.0 * correct / total

    # ----- helpers ------------------------------------------------------------

    def get_dataset_size(self):
        return self.dataset_size
\end{lstlisting}

\subsection{QPSOServer --- Novel Aggregation}
\begin{lstlisting}[caption={QPSO Server Aggregation (server\_qpso.py)}, label={lst:qpso}]
"""
QPSOServer — Quantum-behaved Particle Swarm Optimisation for FL aggregation.

Each client is a "particle".  The server maintains:
  • personal best (pbest)  — best state_dict each client has ever produced.
  • global  best (gbest)   — best state_dict across all clients.
  • mean    best (mbest)   — element-wise mean of all pbests.

Position update (per parameter tensor):
    φ  ~ U(0,1)    u ~ U(0,1)
    p  = φ·mean(pbests) + (1−φ)·gbest          # attraction point
    x' = p  ±  β · |mbest − x| · ln(1/u)       # QPSO step

β (contraction-expansion coefficient) controls exploration vs exploitation.
"""

import copy
import torch
import torch.nn as nn


class QPSOServer:
    """Central server using QPSO-based aggregation."""

    def __init__(self, global_model, clients, device="cuda", beta=0.7):
        self.global_model = global_model
        self.clients      = clients
        self.device       = device
        self.beta         = beta

        # per-client personal bests
        self.personal_best        = {}          # client_id → state_dict
        self.personal_best_scores = {}          # client_id → float (val acc)

        # global best
        self.global_best       = None           # state_dict
        self.global_best_score = 0.0

        # mean best (computed each round)
        self.mean_best = None

        print(f"QPSO Server  |  clients={len(clients)}  β={beta}")

    # ----- initialisation -----------------------------------------------------

    def initialize_particles(self):
        """Set every pbest and gbest to the current global model."""
        state = copy.deepcopy(self.global_model.state_dict())
        for c in self.clients:
            self.personal_best[c.client_id]        = copy.deepcopy(state)
            self.personal_best_scores[c.client_id] = 0.0
        self.global_best       = copy.deepcopy(state)
        self.global_best_score = 0.0
        print("✅ QPSO particles initialised")

    # ----- pbest / gbest updates ----------------------------------------------

    def update_personal_best(self, client_id, weights, val_acc):
        if val_acc > self.personal_best_scores[client_id]:
            self.personal_best[client_id]        = copy.deepcopy(weights)
            self.personal_best_scores[client_id] = val_acc
            return True
        return False

    def update_global_best(self, client_id, val_acc):
        if val_acc > self.global_best_score:
            self.global_best       = copy.deepcopy(self.personal_best[client_id])
            self.global_best_score = val_acc
            return True
        return False

    # ----- mbest ---------------------------------------------------------------

    def calculate_mean_best(self):
        """Element-wise mean of all personal-best state_dicts."""
        first_id = self.clients[0].client_id
        self.mean_best = copy.deepcopy(self.personal_best[first_id])

        for k in self.mean_best:
            self.mean_best[k] = torch.zeros_like(self.mean_best[k],
                                                 dtype=torch.float32)
        for c in self.clients:
            pb = self.personal_best[c.client_id]
            for k in self.mean_best:
                self.mean_best[k] += pb[k].float()

        n = len(self.clients)
        for k in self.mean_best:
            self.mean_best[k] /= n

    # ----- QPSO aggregation ---------------------------------------------------

    def qpso_aggregate(self, client_weights_list):
        """
        Perform one QPSO aggregation step.

        Parameters
        ----------
        client_weights_list : list[(client_id, state_dict, val_acc)]

        Returns
        -------
        aggregated state_dict
        """
        # 1. update personal / global bests
        for cid, w, acc in client_weights_list:
            pb = self.update_personal_best(cid, w, acc)
            gb = self.update_global_best(cid, acc)
            if pb:
                print(f"  {cid}: pbest ↑ {acc:.2f}%")
            if gb:
                print(f"  {cid}: gbest ↑ {acc:.2f}%")

        # 2. mean best
        self.calculate_mean_best()

        # 3. QPSO position update for every parameter key
        agg = copy.deepcopy(self.global_best)

        for k in agg:
            # sum of personal bests  (used to get mean pbest for attraction)
            pbest_sum = torch.zeros_like(agg[k], dtype=torch.float32)
            for c in self.clients:
                pbest_sum += self.personal_best[c.client_id][k].float()

            phi  = torch.rand_like(agg[k].float())
            u    = torch.rand_like(agg[k].float())

            # attraction point
            p = phi * (pbest_sum / len(self.clients)) \
              + (1 - phi) * self.global_best[k].float()

            # sign ±1
            sign = torch.where(
                torch.rand_like(agg[k].float()) < 0.5,
                torch.ones_like(agg[k].float()),
               -torch.ones_like(agg[k].float()),
            )

            mbest   = self.mean_best[k].float()
            current = agg[k].float()

            # QPSO step
            new_val = p + sign * self.beta \
                      * torch.abs(mbest - current) \
                      * torch.log(1.0 / (u + 1e-8))

            agg[k] = new_val.to(self.global_best[k].dtype)

        return agg

    # ----- evaluation ---------------------------------------------------------

    def evaluate_global_model(self, test_loader):
        """Return (accuracy%, loss) on a test DataLoader."""
        self.global_model.eval()
        criterion = nn.CrossEntropyLoss()
        total_loss = 0.0
        correct = total = 0

        with torch.no_grad():
            for images, labels in test_loader:
                images, labels = images.to(self.device), labels.to(self.device)
                outputs = self.global_model(images)
                total_loss += criterion(outputs, labels).item()
                _, pred = outputs.max(1)
                total   += labels.size(0)
                correct += pred.eq(labels).sum().item()

        return 100.0 * correct / total, total_loss / len(test_loader)

    # ----- helpers ------------------------------------------------------------

    def get_global_model(self):
        return copy.deepcopy(self.global_model)
\end{lstlisting}

\subsection{FedAvgServer --- Baseline Aggregation}
\begin{lstlisting}[caption={FedAvg Server (server\_fedavg.py)}, label={lst:fedavg}]
"""
FedAvgServer — Federated Averaging aggregation.

Aggregation rule
    w_global = Σ (n_k / n_total) · w_k
where n_k is the training-set size of client k.
"""

import copy
import torch
import torch.nn as nn


class FedAvgServer:
    """Central server implementing the FedAvg algorithm."""

    def __init__(self, global_model, clients, device="cuda"):
        self.global_model = global_model
        self.clients      = clients
        self.device       = device
        self.total_samples = sum(c.get_dataset_size() for c in clients)

        print(f"FedAvg Server  |  clients={len(clients)}  "
              f"total_samples={self.total_samples}")

    # ----- aggregation --------------------------------------------------------

    def aggregate_weights(self, client_weights):
        """
        Weighted average of client state_dicts.

        Parameters
        ----------
        client_weights : list[(state_dict, dataset_size)]

        Returns
        -------
        aggregated state_dict
        """
        first_sd = client_weights[0][0]
        agg = {}

        for k in first_sd:
            if first_sd[k].is_floating_point():
                # Weighted average for float params
                accumulated = torch.zeros_like(first_sd[k], dtype=torch.float32)
                for state_dict, n_k in client_weights:
                    w = n_k / self.total_samples
                    accumulated = accumulated + state_dict[k].float() * w
                agg[k] = accumulated.to(first_sd[k].dtype)
            else:
                # Non-float (e.g. num_batches_tracked): copy from first client
                agg[k] = first_sd[k].clone()

        return agg

    # ----- evaluation ---------------------------------------------------------

    def evaluate_global_model(self, test_loader):
        """Return (accuracy%, loss) on a test DataLoader."""
        self.global_model.eval()
        criterion = nn.CrossEntropyLoss()
        total_loss = 0.0
        correct = total = 0

        with torch.no_grad():
            for images, labels in test_loader:
                images, labels = images.to(self.device), labels.to(self.device)
                outputs = self.global_model(images)
                total_loss += criterion(outputs, labels).item()
                _, pred = outputs.max(1)
                total   += labels.size(0)
                correct += pred.eq(labels).sum().item()

        return 100.0 * correct / total, total_loss / len(test_loader)

    # ----- helpers ------------------------------------------------------------

    def get_global_model(self):
        return copy.deepcopy(self.global_model)
\end{lstlisting}

\subsection{ResidualLSTM --- Progression Model}
\begin{lstlisting}[caption={LSTM Hybrid Residual Model (infrastructure\_lstm.py, lines 273--349)}, label={lst:lstm}]
class ResidualLSTM(nn.Module):
    """
    LSTM model to predict next residual from previous residuals.
    
    Architecture:
      - Input: sequence of residuals (lookback length)
      - LSTM: capture temporal dependencies in residuals
      - Fully connected layers: output prediction
      - Output: predicted next residual
    """
    
    def __init__(
        self,
        lookback: int = 3,
        hidden_size: int = 32,
        num_layers: int = 1,
        dropout: float = 0.2,
    ):
        super().__init__()
        
        self.lookback = lookback
        self.hidden_size = hidden_size
        self.num_layers = num_layers
        
        # LSTM layers
        self.lstm = nn.LSTM(
            input_size=1,
            hidden_size=hidden_size,
            num_layers=num_layers,
            batch_first=True,
            dropout=dropout if num_layers > 1 else 0,
        )
        
        # Attention layer (optional): weight different timesteps
        self.attention = nn.MultiheadAttention(
            embed_dim=hidden_size,
            num_heads=4,
            dropout=dropout,
            batch_first=True,
        )
        
        # Fully connected layers
        self.fc = nn.Sequential(
            nn.Linear(hidden_size, 64),
            nn.ReLU(),
            nn.Dropout(dropout),
            nn.Linear(64, 32),
            nn.ReLU(),
            nn.Dropout(dropout),
            nn.Linear(32, 1),
        )
    
    def forward(self, x: torch.Tensor) -> torch.Tensor:
        """
        Args:
            x: Tensor of shape (batch_size, lookback)
        
        Returns:
            output: Tensor of shape (batch_size, 1) - predicted residual
        """
        
        # Reshape for LSTM: (batch_size, lookback, 1)
        x = x.unsqueeze(-1)
        
        # LSTM forward
        lstm_out, (h_n, c_n) = self.lstm(x)  # lstm_out: (batch, lookback, hidden)
        
        # Attention: weight the LSTM outputs
        attn_out, _ = self.attention(lstm_out, lstm_out, lstm_out)  # (batch, lookback, hidden)
        
        # Take the final hidden state (could also use attention)
        final_hidden = attn_out[:, -1, :]  # (batch, hidden)
        
        # Fully connected layers
        output = self.fc(final_hidden)  # (batch, 1)
        
        return output.squeeze(-1)  # (batch,)
\end{lstlisting}
